\chapter{Derivation AB: Cluster Expansion of the Logarithm}
\label{ch:derivation_ab_cluster_expansion}

\section{Existence of the Thermodynamic Limit}

To rigorously define the free energy density $f(\beta)$ and prove it is analytic in the strong coupling regime, we employ the Cluster Expansion of the logarithm of the partition function. This derivation (Derivation AB) proves the existence of the thermodynamic limit.

\section{The Mayer Expansion}

For a Lattice Gauge Theory, the partition function is $Z_\Lambda = \int \prod_{b \in \Lambda} dU_b \prod_{p} e^{\frac{\beta}{N} \text{Re Tr} U_p}$.
We expand the Boltzmann weight on each plaquette:
\begin{equation}
e^{\frac{\beta}{N} \text{Re Tr} U_p} = 1 + \rho_p(U_p)
\end{equation}
where $\rho_p$ is small when $\beta$ is small.
Then:
\begin{equation}
Z_\Lambda = \int d\mu \prod_p (1 + \rho_p) = \int d\mu \sum_{P \subset \Lambda} \prod_{p \in P} \rho_p = \sum_{P} W(P)
\end{equation}
where $W(P)$ is the weight of a polymer (set of plaquettes) $P$.

\section{The Logarithm and Connected Clusters}

The logarithm of $Z$ involves a sum over \textit{connected} polymers (clusters) $C$:
\begin{equation}
\ln Z_\Lambda = \sum_{C \subset \Lambda} \Psi(C) W(C)
\end{equation}
where $\Psi(C)$ is a combinatorial factor (Ursell function) derived from the Mobius inversion on the interaction lattice.

\subsection{Convergence via Kotecký-Preiss}
A sufficient condition for the absolute convergence of the series for the free energy density $f_\Lambda = \frac{1}{|\Lambda|} \ln Z_\Lambda$ is:
\begin{equation}
\sum_{C \ni p} |W(C)| e^{a |C|} \le 1
\end{equation}
For small $\beta$, $|W(C)| \approx \beta^{|C|}$. Since the number of clusters of size $k$ grows geometrically (with a constant $\mu_{lattice}$), there exists a radius of convergence $\beta_0$.
For $\beta < \beta_0$, the limit $\Lambda \to \infty$ exists and $f(\beta)$ is analytic.

\chapter{Derivation AD: Duhamel's Expansion}
\label{ch:derivation_ad_duhamel}

\section{Semigroup Perturbation Theory}

To analyze the stability of the mass gap under small perturbations (e.g., adding fermions or changing the action slightly), we compare the semigroups of the perturbed Hamiltonian $H = H_0 + V$.

\section{The Expansion}

The Duhamel formula (or Dyson expansion in imaginary time) expresses the difference:
\begin{equation}
e^{-tH} = e^{-tH_0} - \int_0^t ds e^{-(t-s)H} V e^{-sH_0}
\end{equation}
Iterating this yields:
\begin{equation}
e^{-tH} = e^{-tH_0} \sum_{n=0}^\infty (-1)^n \int_{0 \le s_1 \le \dots \le s_n \le t} ds_1 \dots ds_n V(s_n) \dots V(s_1)
\end{equation}
where $V(s) = e^{sH_0} V e^{-sH_0}$.

\section{Application to Spectral Gaps}
We use this to prove that if $H_0$ has a gap $\Delta_0$ and $V$ is relatively bounded with sufficiently small norm, the gap $\Delta$ of $H$ satisfies:
\begin{equation}
|\Delta - \Delta_0| \le \|V\|
\end{equation}
Crucially, in the LSI framework, we use a hypercontractive version of this expansion to show that the log-Sobolev constant is stable under bounded perturbations of the potential (Holley-Stroock lemma relies on this structure).

\section{Conclusion}
These two derivations provide the "bookends" of the stability analysis: Derivation AB handles the vacuum energy (extensive), while Derivation AD handles the excitation spectrum (intensive).
